% !TeX root = ../papers/nc-p3.tex
\section{一つの隣接対称性から四平方の場合を認識する条件}
\label{sec:adjacent-four-squares}
指定した隣接正証明書で最初の二因子を同一視し、
$w\in\operatorname{Sym}^2 V\otimes W\otimes U$とする。
第三切断の階数6の分解を$P(x):F\to W$, $Q(x):U^*\to F$,
$\dim F=6$と書く。ここで$x\in V^*$であり、
$P(x)Q(y)=P(y)Q(x)$は$w$の最初の二因子を$x,y$で縮約した写像である。
階数分解より$\bigcap_x\ker P(x)=0$、$\sum_x\operatorname{im}Q(x)=F$である。
$w$の対称二次消去子$L$が独立な四つの平方$x_0^2,\ldots,x_3^2$で張られる場合を考える。
このとき$[x_i]$は$\ell^2\in L$を満たす射影点の全体である。実際、混合項の全消滅より
$\ell$の非零座標は一つである。$P_i=P(x_i)$, $Q_i=Q(x_i)$,
$M_{ij}=P_iQ_j$と書くと、$M_{ii}=0$, $M_{ij}=M_{ji}$である。


\begin{proposition}[四本の核直線による対称認識]\label{prop:square-kernel-recognition}
一因子階数を全て4とする。各$i$で$\rk P_i=\rk Q_i=3$、
四本の直線$\ker Q_i$が$U^*$を張るとする。このとき独立な因子の基底変更により
\[
 w_{abcd}=\begin{cases}1&\{a,b,c,d\}=\{0,1,2,3\},\\0&\text{その他}\end{cases}
\]
となる。特に六つの対全てで$\rk M_{ij}=2$ ($i\ne j$)である。
これは四平方の可換完全交叉の対称チャートの一点である。この認識にAS性は仮定しない。
\end{proposition}
\begin{proof}
$u_i\in\ker Q_i$を基底として選び、$K_i=\operatorname{im}Q_i=\ker P_i$、
$I_{ij}=K_i\cap K_j$と置く。$Q_i u_j\in I_{ij}$である。
実際$P_jQ_i u_j=P_iQ_j u_j=0$である。
各$i$で三つの$Q_i u_j$ ($j\ne i$)は$K_i$の基底である。

各無順序対から一本ずつ$v_{ij}=v_{ji}\in I_{ij}$を選ぶ。
各頂点$i$で三つの$v_{ij}$が$K_i$の基底になるという条件は
$\prod_{i<j}I_{ij}$の非空開条件である。固定した$i$については
$v_{ij}=Q_i u_j$と選べるため非空性が従う。
パラメータ空間は既約なので四つの条件を同時に満たせる。
六つの$v_{ij}$は全ての$K_i$を張り、階数分解の全射性よりその和は六次元の$F$である。
よって六つは基底である。従って$I_{ij}=k v_{ij}$であり、
$\rk M_{ij}=\dim K_j-\dim I_{ij}=2$を得る。

非零の$Q_i u_j$と$Q_j u_i$は同じ直線にあるから、
$Q_i u_j=c_{ij}Q_j u_i$、$c_{ji}=c_{ij}^{-1}$と書ける。
相異なる$i,j,k$について
\[
 c_{ik}M_{jk}u_i=M_{ij}u_k=c_{jk}c_{ij}M_{jk}u_i.
\]
$\ker M_{jk}=\langle u_j,u_k\rangle$なので$M_{jk}u_i\ne0$であり、
$c_{ik}=c_{ij}c_{jk}$を得る。各$u_i$をスカラー調整して全ての$c_{ij}=1$にできる。
従って$M_{ij}u_k$は三添字に対称であり、重複添字では零となる。
相異なる三添字から得られる四つのベクトルは一因子階数4より$W$を張る。
欠けた添字で番号付けし、これらを$W$の基底に取れば表示した全係数を得る。
\end{proof}

四平方の対称テンソルでは、$F$を無順序二添字の六次元空間に取り、
$Q_i u_j$を$i\ne j$でその対の基底、$Q_i u_i=0$とすれば、全ての仮定を満たす。
$P_iQ_i=0$の下で両階数が3以上という開条件は両方3と同値であり、
その上で核直線の独立性も開条件である。階数分解の変更は$F$の基底変更であり、
四平方の方向の並べ替えとスカラーは判定を変えない。
従って指定した隣接対称性と四平方の消去子を持つ候補内で、
この非空な相対開部分はただ一つの添字保持同型類を与える。
双対化して$(P_i,Q_i)$を$(Q_i^{\mathsf T},P_i^{\mathsf T})$に替えても命題を適用できる。
よってこの四平方の枝で標準対称類の外に残るには、
どこかの軸で階数対が$(3,3)$と異なるか、全軸が$(3,3)$なら
入力側の四核直線も出力側の四余核の双対直線も、ともに従属でなければならない。
